% Páginas de título (extratitle, frontispiece, titlehead, subject,
% dedication, uppertitleback, lowertitleback, publishers) para el flujo del
% pipeline.
%
% Los comandos KOMA originales solo GUARDAN el contenido (\gdef): el
% renderizado vive en el \maketitle original de KOMA. Como 19-maketitle.tex
% lo redefine, este archivo re-expone los comandos (idénticos a KOMA) y
% define \titlebackformat, el formato que 19-maketitle.tex aplica al
% imprimir el contenido en sus páginas.
%
% El contenido llega serializado a LaTeX por el filter latex/10-titlepages.lua
% (markdown → latex: el doble espacio al final de línea → \\, una línea con
% solo :: → \vspace{\baselineskip}).
\makeatletter
\renewcommand{\extratitle}[1]{\gdef\@extratitle{%
    #1%
}}%
\renewcommand{\frontispiece}[1]{\gdef\@frontispiece{%
    #1%
}}%
\renewcommand{\titlehead}[1]{\gdef\@titlehead{%
    #1%
}}%
\renewcommand{\dedication}[1]{\gdef\@dedication{%
    #1%
}}%
% KOMA define \subject con \newcommand* (no-long): una línea en blanco en el
% argumento rompería la compilación (mismo problema que \subtitle). Long.
\renewcommand{\subject}[1]{\gdef\@subject{%
    #1%
}}%
\renewcommand{\publishers}[1]{\gdef\@publishers{%
    #1%
}}%
\renewcommand{\uppertitleback}[1]{\gdef\@uppertitleback{%
    #1%
}}%
\renewcommand{\lowertitleback}[1]{\gdef\@lowertitleback{%
    #1%
}}%
\makeatother

% Formato del contenido de las páginas internas: footnotesize, interlineado
% 1 y sin indentación en ningún párrafo.
\newcommand{\titlebackformat}{%
  \footnotesize
  \setstretch{1}%
  \setlength{\parindent}{0pt}%
}

% Saltos de página propios para las páginas internas. NO se usan
% \next@tpage/\next@tdpage de KOMA: ejecutan \setparsizes{0}{0} que deja
% \parindent a 0 de forma GLOBAL (el body perdería su indentación). Estos
% solo hacen \clearpage + \thispagestyle{empty} (+ página impar en twoside).
\makeatletter
\newcommand{\titlepage@next}{%
  \clearpage
  \thispagestyle{empty}%
}
\newcommand{\titlepage@nextdouble}{%
  \titlepage@next
  \if@twoside\ifodd\value{page}\else\null\titlepage@next\fi\fi
}
\makeatother

% Colofón final: se imprime al final del documento (el template lo emite
% después del body) y SIEMPRE ocupa solo una página par, la última del
% documento. \titlepage@lasteven es el espejo de \titlepage@nextdouble
% (par → impar): tras \clearpage, si la página actual es impar se inserta
% una página en blanco y el colofón cae en la siguiente par. La página del
% colofón y la insertada no muestran número de página.
% Todo dentro de makeatletter: los reemplazos usan comandos con @
% (\@colophon, \titlepage@lasteven) y fuera de él se tokenizarían como
% \@ (spacefactor) y texto suelto.
\makeatletter
\newcommand{\colophon}[1]{\gdef\@colophon{%
    #1%
}}%
\newcommand{\titlepage@lasteven}{%
  \clearpage
  \thispagestyle{empty}%
  \if@twoside\ifodd\value{page}\null\titlepage@next\fi\fi
}
\newcommand{\colophonpage}{%
  \titlepage@lasteven
  \vspace*{7\baselineskip}%
  \titlepageblock
    {\dimexpr(\linewidth-\colophonwidth)/2\relax}%
    {\dimexpr(\linewidth-\colophonwidth)/2\relax}%
    {\centering\titlebackformat}%
    {\@colophon}%
}
\makeatother
